\section{Discussion}
\label{sec:conclusion}
\arnab{didn't check anything. will depend on the concrete framing and findings.}

We find that code LLMs carry a compact representation of their upcoming sub-goal before writing its implementation. The representation can be extracted from attention-head contributions, verbalized in natural language, and transplanted to trigger the same operation in another programming language. The heads also mediate the model's own implementation: preserving their access to task information largely preserves performance, and restricting a broader set sharply reduces it. These findings connect an interpretable description of what the model is about to do with an internal representation that can influence its behavior.

\paragraph{An abstraction of the upcoming operation.}
The portability of the encoding suggests a separation between the identity of a sub-goal and its realization in code. An operation extracted from Python can guide a C++ implementation, even though the two require different tokens. Natural-language verbalization provides a complementary way to inspect this semantic content. Our evidence concerns the upcoming operation near the start of its implementation; it leaves open how much of a longer plan is represented and how its details are chosen. Verbalizability also does not establish that the representation's native format is natural language.

\paragraph{From goal representation to goal selection.}
\label{sec:mechanism:competence}
The intervention results identify several parts of goal-directed computation. Queries carry information that can favor a particular sub-goal; task-span keys and values supply information to which the model can adapt; and the heads' projections expose the upcoming operation in intermediate residual states. The ablations further distinguish a small set sufficient for readout from a broader set important for implementation. These observations constrain the mechanism, while leaving open how the model combines completed work with the specification to select its next step. Independently varying progress, task order, and target identity offers a way to investigate that computation. The current interventions establish causal involvement under specific access restrictions, rather than a unique bottleneck or an isolated goal-selection circuit.

\paragraph{Scope and limitations.}
Our principal quantitative setting uses explicit task specifications and a known frontier. The oracle measures nontrivial semantic agreement and can credit partial implementations, so success rates should be interpreted separately from executable program correctness. Capture positions are repeated observations of the same specifications, and head-budget comparisons use selections that are not strictly nested. Replications use each model's own heads and activations; differences between models combine checkpoint, head-selection, and generation effects. The few-shot example is qualitative, and evaluating free-form code requires further examples and decisions about sub-goal granularity. These settings are promising tests of how far the code-localized representation generalizes.

\goalscope provides a concrete tool for examining what a model is about to do as generation proceeds. Further work can use these readouts to study how sub-goals change, when an intended operation fails to be implemented, and how the model decides to move on to the next step. Our exploratory boundary-signal interventions (\Cref{app:boundary}) provide a complementary starting point for the last question.
